\documentclass[12pt]{article}
\usepackage[utf8]{inputenc}
\usepackage[T1]{fontenc}
\usepackage{hyperref}
\title{Flinders Island Language Resources for K-12 Curriculum Development}
\author{Compiled by Claude Sonnet 5.5 (claude-sonnet-5-5) for Andr\'as Kornai}
\date{}
\begin{document}
\maketitle

\section{Introduction}
The code fln refers to the extinct Pama-Nyungan Flinders Island language of Cape York, Queensland, not to the Tasmanian Aboriginal varieties associated with Flinders Island in Bass Strait. Wikipedia notes it is unconfirmed as a distinct language. Search results show almost nothing usable for a school curriculum, and nothing here should be taken as a revitalization programme.

\section{Primary Educational Resources}
None found.

\section{Linguistic and Reference Materials}
\begin{itemize}
\item Wikipedia article on the language: \url{https://en.wikipedia.org/wiki/Flinders_Island_language}
\item An ANU Press chapter (PDF) records that Johnny Flinders, one of the last known speakers, was Peter Sutton's principal teacher of his language in 1970--1979: \url{https://press-files.anu.edu.au/downloads/press/n14124/pdf/ch12.pdf}
\end{itemize}

\section{Religious and Community Resources}
None found.

\section{Digital and Audio Resources}
None found.

\section{Curriculum Development Context}
No materials suitable for classroom use were found. Any heritage curriculum would have to start from archival fieldnotes held outside what was surfaced here, in consultation with the relevant Cape York communities.
\end{document}